\documentclass[10pt,a4paper]{amsart}
\usepackage[margin=19mm]{geometry}
\usepackage{amsmath,amssymb,mathtools}
\usepackage[scr=rsfs,cal=euler]{mathalfa}
\usepackage{xfrac}
\usepackage[unicode,hidelinks,bookmarks=false]{hyperref}
\newcommand{\ie}{i.e.\ }
\newcommand{\cf}{cf.\ }
\newcommand{\resp}{respectively\ }
\newcommand{\Subsection}{\subsection}
\providecommand{\nobreakdash}{\nobreak\hbox{-}}
\setlength{\emergencystretch}{2em}
\multlinegap0pt
\usepackage{bm}
\usepackage{mathtools}
\usepackage{url}
\usepackage{color}
\usepackage{enumitem}
\usepackage{booktabs}
\usepackage{placeins}
\newtheorem{dfn}{Definition}[section]
\newtheorem{prop}[dfn]{Proposition}
\newtheorem{prob}[dfn]{Problem}
\newtheorem{lem}[dfn]{Lemma}
\newtheorem{exa}[dfn]{Example}
\newtheorem{thm}[dfn]{Theorem}
\newtheorem{cor}[dfn]{Corollary}
\newtheorem{fact}[dfn]{Fact}
\newtheorem{conj}[dfn]{Conjecture}
\newtheorem{hyp}[dfn]{Hypothesis}
\newtheorem{que}[dfn]{Question}
\newtheorem{rem}[dfn]{Remark}
\newcommand{\1}{
    \bm{1}
}
\newcommand{\pZ}[1]{
(\mathbb{Z}/{ #1 \mathbb{Z}})^{\times}
}
\newcommand{\aZ}[1]{
\mathbb{Z}/{ #1 \mathbb{Z}}
}
\begin{document}
\title[On the level of distribution of Goldbach primes and its appl]{On the level of distribution of Goldbach primes and its applications}
\author{Mizuki Akeno}
\date{}
\maketitle
\noindent\emph{Source and licence.} On the level of distribution of Goldbach primes and its applications. Version arXiv:2606.29559v1 (CC BY 4.0). This material is distributed under the Creative Commons Attribution 4.0 International licence, http://creativecommons.org/licenses/by/4.0/, as recorded in the arXiv OAI licence metadata for this exact version and shown on the abstract page https://arxiv.org/abs/2606.29559v1. Full rights notice: licenses/arxiv-2606.29559-CC-BY-4.0.txt.\par\smallskip
\noindent\textit{Editorial scope.} Counted unit: the original \texttt{thm} environment \texttt{T1p} at source lines 344--360 together with its complete proof at lines 365--376; the delivered document keeps source lines 140--377 so that every referenced label and every citation resolves inside it. Native editable source is the paper's own concatenated TeX; the AMS wrapper is editorial formatting only. No publisher class, upstream script or unrelated artwork is redistributed.
\par\smallskip\noindent\textit{Reference note.} This article ships no compiled bibliography (biblatex); the entries delivered here are composed from the author's own BibTeX (.bib) fields, with no field invented and no entry dropped.
\begin{thm}\label{T1}
Let $\theta \in (0,1/6)$, let $A>0$ and let $k \geq 0$. 
Let $w_1,w_2$ be smooth functions supported on $[\eta,1]$ for some $\eta>0$. 

For all integers $N \in (X/2,X]$ with $O(X(\ln{X})^{-A})$ exceptions, we have
\[ \sum_{d \leq X^{\theta}} \tau(d)^k \sup_{l:(d,l(N-l))=1} |R_{d,l}(N)| \ll \frac{X}{(\ln{X})^A} \]
where
\[ R_{d,l}(N) = \sum_{\substack{m+n=N \\ m \equiv l \bmod d}} \Lambda(m) \Lambda(n) w_1(m/X) w_2(n/X) - \frac{\mathfrak{S}(N)}{\psi_N(d)} \sum_{m+n=N} w_1(m/X) w_2(n/X) \]
and $\psi_N$ is defined by
\[ \psi_N(d) = d \prod_{\substack{p|d \\ p | N}} \left( 1 - \frac{1}{p} \right) \prod_{\substack{p|d \\ p \nmid N}} \left( 1 - \frac{2}{p} \right). \]
Here and throughout, we use the convention that the supremum over an empty set is $0$. 
\end{thm}

Indeed, we prove a more general version in which one of the $\Lambda$ factors is replaced by a more general sequence such as the indicator function of rough numbers. See Theorem \ref{T1p}. 


By a slight modification of the result of Maier-Pomerance \cite[Theorem 3.1]{maier1990unusually} concerning the distribution of generalized twin primes in arithmetic progressions, one can establish such a result with $\theta < c$ for some absolute constant $c>0$. Their argument is based on the work of Montgomery and Vaughan \cite{montgomery1975exceptional} on the bounds for the exceptional set in the Goldbach conjecture. 
A computation by the author using standard zero-density estimates for Dirichlet $L$-functions shows that one may take $c = \frac{10}{314} = 0.0318\ldots$, although this value is probably not optimal. 


We use methods that are different from those of \cite{maier1990unusually} in both major and minor arcs. 
We establish a new exponential sum estimate. See Theorem \ref{Lmb}. The exponent $1/6$ in Theorem \ref{T1} comes from this estimate. 
On the major arc, we use a mean value theorem of Choi and Kumchev \cite{choi2006mean} for Dirichlet polynomials, together with an argument of Maynard \cite{maynard2024note}. 


As an application of Theorems \ref{T1} and \ref{T1p}, we consider the solvability of the equations of the form $p_1+p_2=N, f(p_1,p_2) \in \mathbb{P}_{r}$ for some function $f:\mathbb{Z} \times \mathbb{Z} \to \mathbb{Z}$ and $r \geq 2$. 



Matom{\"a}ki and Z{\'u}{\~n}iga-Alterman \cite{matomaki2025weighted} showed that one can find $\mathbb{P}_3$ if the original and the switched problem have level of distribution at least $0.267$. 
We extend their result and show that one can find $\mathbb{P}_4$ if the original and the switched problem have level of distribution at least $0.1635$. See also the Appendix for similar results for $\mathbb{P}_k, 2 \leq k \leq 7$. 
Since $1/6=0.166\ldots>0.1635$, we obtain the following result. 
\begin{thm}\label{T2}
Let $A>0$ and $X \geq 2$. 
For all but $O(X(\ln{X})^{-A})$ even integers $N \in (X/2,X]$, there exist primes $p_1,p_2$ such that $p_1+p_2=N$ and $p_1-p_2+1$ has at most $4$ prime factors. 
\end{thm}
The numerical comparison in the Appendix indicates that the weights appearing in previous applications of the weighted sieve do not suffice for this purpose, since the resulting criteria remain above the level $1/6$. See the Appendix.

We also use Theorem \ref{T1} and Richert's weighted sieve to obtain
\begin{thm}\label{T3}
Let $A>0$ and $X \geq 2$. 
For all but $O(X(\ln{X})^{-A})$ integers $N \in (X/2,X]$ with $6|N$, there exist primes $p_1,p_2$ such that $p_1+p_2=N$ and $2p_1p_2+1$ has at most $13$ prime factors. 
\end{thm}
Note that the number $13$ arises as $13=2 \times 6 + 1$.






\subsection{AI tool disclosure}


OpenAI Codex (GPT-5.5) was used in the preparation of this manuscript in the following ways:
\begin{itemize}
\item assistance with the numerical computations reported in the Appendix and with the proof of the asymptotic statement there in the regime \(k\to\infty\);
\item assistance with searching for existing results relevant to the manuscript;
\item proofreading for minor typographical, stylistic, and mathematical errors.
\end{itemize}
Apart from these uses, the arguments and proofs in the paper were written by the author.  All AI-assisted material used in the proofs was checked by the author, who takes full responsibility for the content of the paper.


\section{Exponential sum estimate}


In this section, we prove the following estimate for exponential sums. 


\begin{thm}\label{Lmb}
Let $D \leq X$ and $\varepsilon>0$. 
Let $w:\mathbb{R} \to \mathbb{R}$ be a smooth function supported on $[0,1]$. 
Let $\alpha \in [0,1], q \geq 1, (q,a)=1, |\alpha-a/q| \leq 1/q^2$. Then,
\begin{equation} \sum_{d \sim D} \sum_{b \in \aZ{d}} \left| \sum_{n} \Lambda(n) e((\alpha+b/d)n) w(n/X) \right|^2 \ll X^{2+\varepsilon} \left( \frac{D}{q} + \frac{qD}{X} + D^2 X^{-1/3} \right) \label{EXP} \end{equation}
\end{thm}

For the proof of Theorem \ref{T1}, a weaker form of \eqref{EXP} with $\aZ{d}$ replaced by $\pZ{d}$ would suffice. However, the argument below treats both cases in the same way, so we state the estimate in the present form.
In this setting, the large sieve gives the bound $\ll X^2(\ln X)^2$ when $D \leq X^{1/2}$, whereas Theorem \ref{Lmb} gives a sharper estimate when $q \in [X^{\varepsilon}D,X^{1-\varepsilon}D^{-1}]$ and $D \leq X^{1/6 - \varepsilon}$ for some $\varepsilon>0$.  



The following two lemmas are fundamental. 

For $x \in \mathbb{R}$, let $\|x\|$ denote the distance from $x$ to the nearest integer. 
\begin{lem}\label{sumw}
Let $w:\mathbb{R} \to \mathbb{R}$ be a smooth, compactly supported function. 
Let $\alpha \in \mathbb{R}, \varepsilon>0$. Then
\begin{equation} \sum_{n} e(\alpha n) w(n/X) \ll X\#\{ n^* \in \mathbb{Z} : |\alpha - n^*| \leq X^{\varepsilon-1} \} + X^{-100} \label{ew_1} \end{equation}
and
\begin{equation} \sum_{n} e(\alpha n) w(n/X) \ll \frac{1}{\|\alpha\|}. \label{ew_2} \end{equation}
Here, the implicit constant depends only on $\varepsilon$ and $\{\|w^{(j)}\|_{\infty} : j \leq \lceil 200/\varepsilon \rceil \}$. 
\end{lem}

\begin{proof}
By the Poisson summation formula, we have
\[ \sum_{n} e(\alpha n) w(n/X) = X \sum_{n^*} \hat{w}(X(n^*-\alpha)), \]
where $\hat{w}(x) = \int_{\mathbb{R}} w(y) e(-xy) dy$. Integrating by parts $j$ times, we see that
\[ \hat{w}(x) = \frac{1}{(2\pi i x)^j} \int_{\mathbb{R}} w^{(j)}(y) e(-xy) dy \ll_j \|w^{(j)}\|_{\infty} |x|^{-j} \quad (x \neq 0) \]
and therefore
\[ X \hat{w}(X(n^*-\alpha)) \ll \frac{X}{(1+X|\alpha-n^*|)^j} \]
for any $j \geq 0$. Thus, the contribution from the terms with $|\alpha-n^*|>1/2$ is $O(X^{-100})$. 
If $|\alpha-n^*|\leq1/2$, we take $j=1, \lceil 200/\varepsilon \rceil$. This yields \eqref{ew_1} and \eqref{ew_2}. 
\end{proof}

\begin{lem}\label{bohr}
Let $\alpha \in \mathbb{R}, q \geq 1, (a,q)=1, |\alpha-a/q| \leq 1/q^2$ and $M,\delta>0$. 
\begin{equation} \#\{ M<m \leq 2M,n \in \mathbb{Z} : |\alpha - n/m| \leq \delta \} \ll M \left( \frac{1}{q} + q \delta + \delta M \right)  \label{qmn} \end{equation}
\end{lem}
\begin{proof}
If $M \delta \geq 1/4$, then there are $O(M\delta)$ integers $n$ for each $m$, and thus the claim follows. Thus, we may suppose that $M \delta < 1/4$. In this case, the left-hand side of \eqref{qmn} is bounded by
\[ \#\{ M<m \leq 2M : \| \alpha m\| \ll \delta M \} \]
since $n$ is unique for each $m$. Write $m=qs+r$ with $0< r \leq q$. Then $\|\alpha m\| =\|ar/q+(qs+r) \epsilon \|$, where $\epsilon=\alpha-a/q$.  

If $M \leq q/4$, then $s=0$ and $0<r\leq q/2$. Thus, we have
\[ \|\alpha m\| \geq \|ar/q\|-|\epsilon| q/2 \geq \|ar/q\|/2 \]
and therefore, there are $O(q \delta M)$ possible values of $r$. This gives \eqref{qmn}. 

If $M \geq q/4$, there are $O(M/q)$ possible values of $s$, and for each $s$, there are $q \delta M+1$ possible values of $r$. This gives \eqref{qmn}. 
\end{proof}




\begin{lem}\label{bohr_dn}
Let $\varepsilon>0$, and let $D,N \leq X$. Let $\alpha \in [0,1]$ with $q \geq 1, (a,q)=1$ and $|\alpha-a/q| \leq 1/q^2$. Then we have
\begin{equation} \label{exp_I_II} \#\{n \sim N, m \in \mathbb{Z},d \sim D, 1 \leq b \leq d : |\alpha+b/d + m/n| \leq \delta \} \ll X^{\varepsilon} ND \left( \frac{1}{q} + q \delta + \delta ND \right) \end{equation}
for all $\delta>0$. Here, the implicit constant depends only on $\varepsilon$. 
\end{lem}

\begin{proof}
Since $\alpha+b/d \ll 1$, we may suppose that $m \ll n$. 
For each $d,n$, let $g=(d,n)$ and write $d=d'g, n=n'g$. For each $c \ll d'n'g$, we consider the equation
\[ \frac{b}{d} + \frac{m}{n} = \frac{c}{d'n'g} \]
where $(b,m) \in \mathbb{Z}^2, b \ll d, m \ll n$. Since $bn' \equiv c \bmod d'$ and $(d',n')=1$, there are $\ll g$ possible values of $b$. For each $b$, $m$ is unique. Thus the number of solutions to this equation is $\ll g$. Hence, the left-hand side of \eqref{exp_I_II} is
\begin{align*}
&\ll (\ln{X}) \max_{G \ll \min(D,N)} G \#\{n' \sim N/G, d' \sim D/G, g \sim G, c \ll d'n'g : |\alpha+c/d'n'g| \leq \delta \}, \\
&\ll X^{\varepsilon} \max_{G \ll \min(D,N)} G \#\{m \sim ND/G, n \ll m : |\alpha+n/m| \leq \delta \}. 
\end{align*}
This and Lemma \ref{bohr} complete the proof. 
\end{proof}


We now prove Theorem \ref{Lmb}. 
The proof is based on a straightforward application of the above lemmas. 
A variant using large-sieve-type estimates might lead to an improved result. 


\begin{proof}[Proof of Theorem \ref{Lmb}]

For any given $U,V \geq 1$, by Vaughan's identity, $\sum_{n} \Lambda(n) e((\alpha+b/d)n) w(n/X)$ can be written as a linear combination of $O((\ln{X})^2)$ sums of the form
\[ \sum_{\substack{m,n \\ m \sim M}} a_m e((\alpha+b/d)mn) w(mn/X), \]
where $M \leq UV$ and $(a_m)$ is a sequence with $|a_m| \leq X^{o(1)}$, and
\[ \sum_{\substack{m \sim M \\ n \sim N}} a_m b_n e((\alpha+b/d)mn) w(mn/X), \]
where $M \geq U, N \geq V, MN \asymp X$ and $(a_m),(b_n)$ are sequences with $|a_m|,|b_n| \leq X^{o(1)}$. Thus, setting $U=V=X^{1/3}$, the theorem follows from the Type I bound
\begin{equation}
\sum_{d \sim D} \sum_{b \in \aZ{d}} \left| \sum_{\substack{m,n \\ m \sim M}} a_m e((\alpha+b/d)mn) w(mn/X) \right|^2 \ll X^{2+2\varepsilon} \left( \frac{D}{q} + \frac{qD}{X} + \frac{MD^2}{X} \right), \label{exp_I}
\end{equation}
and the Type II bound
\begin{equation}
\sum_{d \sim D} \sum_{b \in \aZ{d}} \left| \sum_{\substack{m \sim M \\ n \sim N}} a_m b_n e((\alpha+b/d)mn) w(mn/X) \right|^2 \ll X^{2+2\varepsilon} \left( \frac{D}{q} + \frac{qD}{X} + \frac{MD^2}{X} + \frac{ND^2}{X} \right), \label{exp_II}
\end{equation}
for any given $\varepsilon>0$, $M,N$ with $MN \asymp X$ and sequences $(a_m),(b_n)$ with $|a_m|,|b_n|\leq 1$. 


We show \eqref{exp_II}. By the Cauchy--Schwarz inequality and interchanging the order of summation, we have
\begin{align*}
\left| \sum_{\substack{m \sim M \\ n \sim N}} a_m b_n e((\alpha+b/d)mn) w(mn/X) \right|^2 &\leq M \sum_{m} \left| \sum_{n \sim N} b_n e((\alpha+b/d)mn) w(mn/X) \right|^2, \\
&\leq M \sum_{n_1,n_2 \sim N} \left| \sum_{m} e((\alpha+b/d)(n_1-n_2)m) w(mn_1/X) w(mn_2/X) \right|. \\
\end{align*}
Since $w(\cdot n_1/N) \times w(\cdot n_2/N)$ is a smooth, compactly supported function and its derivatives are bounded with constants independent of $n_1,n_2$, by Lemma \ref{sumw}, this is
\begin{align*}
&\ll M^2 \sum_{n_1,n_2 \sim N} \#\{m^* \in \mathbb{Z}:|(\alpha+b/d)(n_1-n_2)-m^*| \leq N/X^{1-\varepsilon} \} + X^{-10}, \\
&\ll XM + X M \#\left\{1 \leq n \leq N, m^* \in \mathbb{Z}:|\alpha+b/d - m^*/n| \leq \frac{N}{n X^{1-\varepsilon}} \right\} + X^{-10}.
\end{align*}
We divide $n$ into dyadic ranges, and use Lemma \ref{bohr_dn}. This gives
\begin{align*}
&\sum_{d \sim D} \sum_{b \in \aZ{d}} \#\left\{1 \leq n \leq N, m^* \in \mathbb{Z}:|\alpha+b/d - m^*/n| \leq \frac{N}{n X^{1-\varepsilon}} \right\}, \\
&\ll X^{2\varepsilon} \left( \frac{ND}{q} + \frac{qND}{X} + \frac{(ND)^2}{X} \right).
\end{align*}
This gives \eqref{exp_II}. 


To prove \eqref{exp_I}, we simply apply
\[ \left| \sum_{\substack{m,n \\ m \sim M}} a_m e((\alpha+b/d)mn) w(mn/X) \right| \ll X, \]
and
\[ \left| \sum_{\substack{m,n \\ m \sim M}} a_m e((\alpha+b/d)mn) w(mn/X) \right| \ll \frac{X}{M} \#\{ m \sim M, n^* \in \mathbb{Z} : |\alpha+b/d-n^*/m| \ll X^{\varepsilon-1} \} + X^{-10}. \]
Thus the left-hand side of \eqref{exp_I} is bounded by
\[ \frac{X^2}{M} \#\{ d \sim D, b \in \aZ{d}, m \sim M, n^* \in \mathbb{Z} : |\alpha+b/d-n^*/m| \ll X^{\varepsilon-1} \} + X^{-1}. \]
This and Lemma \ref{bohr_dn} give \eqref{exp_I}. 
\end{proof}







\section{Level of distribution of Goldbach primes}




To prove Theorem \ref{T2}, we need Theorem \ref{T1} and an analogue of Theorem \ref{T1} for equations of the form $2m+n=N+1$, in which one of the $\Lambda$ factors is replaced by a structured weight. 
Since the proofs are essentially the same, we prove a generalized version in both settings. 


\begin{thm}\label{T1p}
Let $\theta \in (0,1/6)$ and let $\eta>0$ be sufficiently small. 
Let $k \geq 1$ and $\beta$ be a function of the form
\[ \beta(n) = \sum_{\substack{n_1 \cdots n_k=n \\ N_i<n_i \leq N'_i}} \Lambda(n_1) \cdots \Lambda(n_k), \]
where $N_1=1,N'_1=X$ if $k=1$, and $X^{\eta}<N_i <N'_i \leq 2N_i$, $N_1 \cdots N_k \leq X$ if $k>1$. 
Let $w_1,w_2$ be smooth functions supported on $[\eta,1]$. 

Let $A>0$ and $X \geq 2$. For all but $O(X(\ln{X})^{-A})$ integers $N \in (X/2,X]$, we have
\[ \sum_{d \leq X^{\theta}} \sup_{l:(d,l(N-l))=1} \left| \sum_{\substack{m+n=N \\ m \equiv l \bmod d}} \Lambda(m) \beta(n) w_1(m/X) w_2(n/X) - M_0 \right| \ll \frac{X}{(\ln{X})^A} \]
where $M_0=M_0(N,d)$ is given by
\[ M_0 = \frac{\mathfrak{S}(N)}{\psi_N(d)} \sum_{m+n=N} \beta(n) w_1(m/X) w_2(n/X), \quad \mathfrak{S}(N) = 2 \1_{2 | N} \prod_{p>2} \left( 1 + \frac{c_p(N)}{\varphi(p)^2} \right). \]

For all but $O(X(\ln{X})^{-A})$ integers $N \in (X/2,X]$, we have
\[ \sum_{\substack{d \leq X^{\theta} \\ 2 \nmid d}} \sup_{l:(d,l(N-2l))=1} \left| \sum_{\substack{2m+n=N \\ m \equiv l \bmod d}} \Lambda(m) \beta(n) w_1(m/X) w_2(n/X) - M'_0 \right| \ll \frac{X}{(\ln{X})^A} \]
where $M'_0=M'_0(N,d)$ is given by
\[ M'_0 = \frac{\mathfrak{S}'(N)}{2\psi_N(d)} \sum_{2m+n=N} \beta(n) w_1(m/X) w_2(n/X), \quad \mathfrak{S}'(N) = 2 \1_{2 \nmid N} \prod_{p>2} \left( 1 + \frac{c_p(N)}{\varphi(p)^2} \right). \]
\end{thm}

\begin{proof}[Proof of Theorem \ref{T1} assuming Theorem \ref{T1p}]
If $k=0$, the claim follows from Theorem \ref{T1p} with $\beta = \Lambda$ and the prime number theorem. 

By the triangle inequality, we have 
\[ R_{d,l}(N) \ll \frac{X}{d} (\ln{X})^2 \]
for all $d,l$ and $N \leq X$. Hence, the Cauchy--Schwarz inequality gives
\begin{align*}
\sum_{d \leq X^{\theta}} \tau(d)^k \sup_{l:(d,l(N-l))=1} |R_{d,l}(N)| &\ll \left( \sum_{d \leq X^{\theta}} \tau(d)^{2k} \frac{X}{d} (\ln{X})^2 \right)^{1/2} \left( \sum_{d \leq X^{\theta}} \sup_{l:(d,l(N-l))=1} |R_{d,l}(N)| \right)^{1/2} \\
&\ll (\ln{X})^{O_k(1)} \left( X \sum_{d \leq X^{\theta}} \sup_{l:(d,l(N-l))=1} |R_{d,l}(N)| \right)^{1/2}. \\
\end{align*}
This reduces the general case to the case $k=0$. 
\end{proof}


\begin{thebibliography}{99}
\bibitem[]{choi2006mean} Choi, Stephen Kwok-Kwong; Kumchev, Angel V.. Mean values of {Dirichlet} polynomials and applications to linear equations with prime variables. Acta Arithmetica 123 (2006) 125--142
\bibitem[]{maier1990unusually} Maier, Helmut; Pomerance, Carl. Unusually large gaps between consecutive primes. Transactions of the American Mathematical Society 322 (1990) 201--237
\bibitem[]{matomaki2025weighted} Matom{\"a}ki, Kaisa; Z{\'u}{\~n}iga-Alterman, Sebastian. Weighted sieves with switching. Mathematical Proceedings of the Cambridge Philosophical Society 179 (2025) 351--372
\bibitem[]{maynard2024note} Maynard, James. A note on integers expressible as the sum of three squares of primes. Acta Arithmetica 214 (2024) 399--419
\bibitem[]{montgomery1975exceptional} Montgomery, H; Vaughan, R. The exceptional set of Goldbach's problem. Acta arithmetica 27 (1975) 353--370
\end{thebibliography}
\end{document}
